\chapter{Planning as an Architectural Control}
\label{ch:planning}

\section{Planning Before the Moment of Action}

\begin{theorem}[Planning theorem]
Preparation performed during a high-clarity period lowers future activation
barriers more reliably than last-minute self-persuasion.
\end{theorem}

\begin{discussion}
Planning works on at least four channels:
\begin{enumerate}[nosep]
  \item it removes ambiguity from the target state;
  \item it shortens the chain of in-the-moment decisions;
  \item it places material cues into the environment;
  \item it can import external force through schedules, reminders, or other
        people.
\end{enumerate}
As a result, planning is not merely motivational language. It is a structural
rewrite of future choice conditions.
\end{discussion}

\section{A Planning Stack}

For teaching purposes, it is useful to separate planning into layers:
\begin{enumerate}[nosep]
  \item \textbf{Outcome layer:} what end state is desired?
  \item \textbf{Session layer:} what concrete session will move that outcome?
  \item \textbf{Entry layer:} what exact first action starts the session?
  \item \textbf{Environment layer:} what physical arrangement lowers entry
        friction?
  \item \textbf{Trigger layer:} when and where is the decision window expected?
\end{enumerate}

\begin{proposition}[Binary execution advantage]
When a plan reduces a future action to a binary choice---execute or refuse---it
consumes less volitional energy than a plan that still requires on-the-spot
design.
\end{proposition}

\section{Chapter-Level Workflow}

\begin{templatebox}
\textbf{Minimal transition plan}
\begin{enumerate}[nosep]
  \item Name the target state.
  \item Name the likely decision window.
  \item Write the first visible action.
  \item Place the required objects in the future environment.
  \item Define the stop rule for the session.
\end{enumerate}
\end{templatebox}

\section{Why Unplanned Intentions Fail}

\begin{corollary}[The unplanned agent follows current inertia]
If no preparation lowers future barrier terms, then the default winner of the
next decision is the present state rather than the verbally preferred state.
\end{corollary}

\begin{example}[Exercise after work]
The sentence ``I should work out after work'' is not yet a usable plan.
It leaves open the location, bag, sequence, exercise selection, time length,
and exit condition. Those unanswered questions are hidden barrier terms. A
real plan resolves them upstream.
\end{example}

\section{Recommended Reading}

\begin{readingbox}
\textbf{Foundation}
\begin{itemize}[nosep]
  \item Stanislas Dehaene, \emph{Consciousness and the Brain}, for the broader
        logic of structured access and global coordination.
\end{itemize}
\textbf{Modern Research}
\begin{itemize}[nosep]
  \item Research on conscious-state signatures and stability helps motivate why
        planning should be treated as state engineering rather than pep talk.
\end{itemize}
\textbf{Frontier}
\begin{itemize}[nosep]
  \item Frontier physical theories are not needed to justify planning as an
        architectural control.
\end{itemize}
\end{readingbox}

\begin{summarybox}
This chapter turns planning into infrastructure. Its value for the project is
that it gives later case studies and exercises a stable template: identify the
window, lower ambiguity, preload the environment, and collapse future choice.
\end{summarybox}
